\section{Renormalization group equation and the asymptotic formula}
\label{sec:4}
\subsection{Renormalization group equation for the coefficients}
\label{sec:4.1}
We now operate
\begin{equation}
   \left(\mu\frac{\partial}{\partial\mu}\right)_0,
\label{eq:(4.1)}
\end{equation}
on both sides of Eq.~\eqref{eq:(3.7)}. Since the left-hand side
of~Eq.~\eqref{eq:(3.7)}, i.e., Eq.~\eqref{eq:(3.5)}, is entirely expressed by
bare quantities through the flow equation~\eqref{eq:(3.1)} and the initial
condition~\eqref{eq:(3.4)}, the action of~\eqref{eq:(4.1)} on the left-hand
side identically vanishes. On the right-hand side, this vanishing must hold in
each power of~$t$. Thus we infer that
\begin{align}
   &\left(\mu\frac{\partial}{\partial\mu}\right)_0c_T(t)
   \left\{T_{\mu\nu}\right\}_R(x)=0,
\label{eq:(4.2)}
\\
   &\left(\mu\frac{\partial}{\partial\mu}\right)_0c_S(t)
   \left\{\frac{1}{4}F_{\rho\sigma}^aF_{\rho\sigma}^a\right\}_R(x)=0.
\label{eq:(4.3)}
\end{align}

For the first relation~\eqref{eq:(4.2)}, we recall that the energy--momentum
tensor is not renormalized as~Eq.~\eqref{eq:(2.6)}. Then, by expressing the
operation~\eqref{eq:(4.1)} in terms of renormalized quantities, we have
\begin{equation}
   \left(\mu\frac{\partial}{\partial\mu}
   +\beta\frac{\partial}{\partial g}\right)c_T(t)=0.
\label{eq:(4.4)}
\end{equation}

For Eq.~\eqref{eq:(4.3)}, on the other hand, from~Eq.~\eqref{eq:(2.11)},
\begin{equation}
   \left(\mu\frac{\partial}{\partial\mu}
   +\beta\frac{\partial}{\partial g}
   +\gamma_S\right)c_S(t)=0,
\label{eq:(4.5)}
\end{equation}
where
\begin{equation}
   \gamma_S
   \equiv-\left(\mu\frac{\partial}{\partial\mu}\right)_0\ln Z_S.
\label{eq:(4.6)}
\end{equation}
Equations~\eqref{eq:(2.21)}, \eqref{eq:(2.7)}, and~\eqref{eq:(2.17)} yield
\begin{equation}
   \gamma_S=-g^3\frac{d}{dg}\left(\frac{\beta}{g^3}\right)
   =2b_1g^4+O(g^6).
\label{eq:(4.7)}
\end{equation}

Similarly, for~Eq.~\eqref{eq:(3.8)}, we have
\begin{align}
   &\left(\mu\frac{\partial}{\partial\mu}
   +\beta\frac{\partial}{\partial g}\right)
   \left\langle E(t,x)\right\rangle=0,
\label{eq:(4.8)}
\\
   &\left(\mu\frac{\partial}{\partial\mu}
   +\beta\frac{\partial}{\partial g}
   +\gamma_S\right)c_E(t)=0,
\label{eq:(4.9)}
\end{align}
and thus
\begin{equation}
   \left(\mu\frac{\partial}{\partial\mu}
   +\beta\frac{\partial}{\partial g}\right)\frac{c_S(t)}{c_E(t)}=0.
\label{eq:(4.10)}
\end{equation}

By the standard argument and from the fact that dimensionless quantities can
depend on the renormalization scale~$\mu$ only through the dimensionless
combination~$\sqrt{8t}\mu$, the above renormalization group equations imply
that
\begin{align}
   &c_T(t)(g;\mu)=c_T(t_0)(\Bar{g}(-\xi);\mu_0),
\label{eq:(4.11)}
\\
   &c_S(t)(g;\mu)
   =\exp\left[\int_0^{-\xi}d\xi'\,\gamma_S\left(\Bar{g}(\xi')\right)\right]
   c_S(t_0)(\Bar{g}(-\xi);\mu_0),
\label{eq:(4.12)}
\\
   &t^2\left\langle E(t,x)\right\rangle(g;\mu)
   =t_0^2\left\langle E(t_0,x)\right\rangle(\Bar{g}(-\xi);\mu_0),
\label{eq:(4.13)}
\\
   &\frac{c_S(t)}{c_E(t)}(g;\mu)=\frac{c_S(t_0)}{c_E(t_0)}(\Bar{g}(-\xi);\mu_0),
\label{eq:(4.14)}
\end{align}
where the dependence on the renormalized gauge coupling and on the
renormalization scale has been explicitly written. In these expressions, the
running coupling~$\Bar{g}(\xi)$ is defined by
\begin{equation}
   \frac{d\Bar{g}(\xi)}{d\xi}=\beta\left(\Bar{g}(\xi)\right),\qquad
   \Bar{g}(0)=g,
\label{eq:(4.15)}
\end{equation}
and we introduce a variable
\begin{equation}
   \xi\equiv\ln\frac{\sqrt{8t}\mu}{\sqrt{8t_0}\mu_0}.
\label{eq:(4.16)}
\end{equation}

In the one-loop order, the running couping~\eqref{eq:(4.15)} is given by
\begin{equation}
   \Bar{g}(-\xi)^2=\frac{1}{2b_0}\frac{1}{-\xi+1/(2b_0g^2)}
   =\frac{1}{2b_0}\frac{1}{-\ln(\sqrt{8t}\Lambda)+\ln(\sqrt{8t_0}\mu_0)},
\label{eq:(4.17)}
\end{equation}
where $\Lambda$ is the $\Lambda$~parameter in the one-loop level,
\begin{equation}
   \Lambda=\mu e^{-1/(2b_0g^2)},
\label{eq:(4.18)}
\end{equation}
and the integral appearing in~Eqs.~\eqref{eq:(4.12)} is
\begin{equation}
   \int_0^{-\xi}d\xi'\,
   \gamma_S\left(\Bar{g}(\xi')\right)
   =\frac{b_1}{b_0}
   \left[
   g^2-\Bar{g}(-\xi)^2
   \right].
\label{eq:(4.19)}
\end{equation}
In the small flow-time limit $t\to0$, $-\xi\to+\infty$ and the running
coupling~$\Bar{g}(-\xi)$~\eqref{eq:(4.17)} becomes very small thanks to the
asymptotic freedom. Thus, the right-hand sides
of~Eqs.~\eqref{eq:(4.11)}--\eqref{eq:(4.14)} allow us to compute the small
flow-time behavior of the coefficients by using the perturbation theory.

\subsection{Lowest-order approximation and the asymptotic formula}
By substituting the solution of the flow equation~\eqref{eq:(3.1)} (see
Ref.~\cite{Luscher:2011bx}) in the tree-level approximation
to~Eq.~\eqref{eq:(3.7)}, we have
\begin{equation}
   c_T(t)=g_0^2,
\label{eq:(4.20)}
\end{equation}
simply because our energy--momentum tensor~\eqref{eq:(2.4)} is proportional
to~$1/g_0^2$. If we apply the right-hand side of~Eq.~\eqref{eq:(4.11)} to this
expression by substituting Eq.~\eqref{eq:(4.17)}, however, it depends
on~$\sqrt{8t_0}\mu_0$ while the left-hand side of~Eq.~\eqref{eq:(4.11)} does
not. This shows that $c_T(t)$ should depend on~$g^2$ and~$\sqrt{8t}\mu$ through
a particular combination as (for $D=4$)
\begin{equation}
   c_T(t)=
   g^2\left\{1+2b_0g^2\left[\ln(\sqrt{8t}\mu)+c_1\right]+O(g^4)\right\},
\label{eq:(4.21)}
\end{equation}
where $c_1$ is a constant. Similarly, since the lowest-order approximation
in~Eqs.~\eqref{eq:(3.7)} and~\eqref{eq:(3.8)} yields
\begin{equation}
   c_E(t)=1,\qquad
   c_S(t)=-\frac{b_0}{2}g_0^2\mu^{-2\epsilon},
\label{eq:(4.22)}
\end{equation}
where the latter follows from~Eq.~\eqref{eq:(3.10)}, from~Eq.~\eqref{eq:(4.14)}
we have
\begin{equation}
   \frac{c_S(t)}{c_E(t)}
   =-\frac{b_0}{2}
   g^2\left\{1+2b_0g^2\left[\ln(\sqrt{8t}\mu)+c_2\right]+O(g^4)\right\},
\label{eq:(4.23)}
\end{equation}
where $c_2$ is another constant.\footnote{Using Eq.~\eqref{eq:(4.21)}
in~Eq.~\eqref{eq:(3.10)}, we have
\begin{equation}
   c_S(t)=-\frac{b_0}{2}
   g^2\left\{
   1+2b_0g^2
   \left[
   \ln(\sqrt{8t}\mu)+c_1+\frac{b_1}{2b_0^2}\right]+O(g^4)
   \right\},
\label{eq:(4.24)}
\end{equation}
and then using Eq.~\eqref{eq:(4.23)},
\begin{equation}
   c_E(t)=1+2b_0g^2
   \left(c_1-c_2+\frac{b_1}{2b_0^2}\right)+O(g^4).
\label{eq:(4.25)}
\end{equation}
}

Applying Eqs.~\eqref{eq:(4.11)} and~\eqref{eq:(4.14)} to above expressions
and using~Eq.~\eqref{eq:(4.17)}, we finally have the asymptotic behaviors of
the coefficients in~Eq.~\eqref{eq:(3.11)},
\begin{equation}
   \frac{1}{c_T(t)}
   \stackrel{t\to0+}{\sim}
   -2b_0\left[\ln(\sqrt{8t}\Lambda)+c_1\right]
\label{eq:(4.26)}
\end{equation}
and
\begin{equation}
   \frac{c_S(t)}{c_E(t)}
   \stackrel{t\to0+}{\sim}
   -\frac{b_0}{2}
   \frac{1}{-2b_0\left[\ln(\sqrt{8t}\Lambda)+c_2\right]},
\label{eq:(4.27)}
\end{equation}
and hence
\begin{equation}
   \frac{c_S(t)}{c_T(t)c_E(t)}
   \stackrel{t\to0+}{\sim}
   -\frac{b_0}{2}
   \left[1-\frac{c_1-c_2}{-\ln(\sqrt{8t}\Lambda)}\right].
\label{eq:(4.28)}
\end{equation}
That is,
\begin{align}
   \left\{T_{\mu\nu}\right\}_R(x)
   &\stackrel{t\to0+}{\sim}
   \biggl\{
   -2b_0\left[\ln(\sqrt{8t}\Lambda)+c_1\right]U_{\mu\nu}(t,x)
\notag\\
   &\qquad\qquad{}
   +\frac{b_0}{2}
   \left[1-\frac{c_1-c_2}{-\ln(\sqrt{8t}\Lambda)}\right]
   \delta_{\mu\nu}
   \left[
   E(t,x)-\left\langle E(t,x)\right\rangle\right]
   \biggr\}.
\label{eq:(4.29)}
\end{align}
This is the relation that we were seeking: One can obtain the
correctly-normalized conserved energy--momentum tensor from the small flow-time
behavior of gauge-invariant products given by the Yang--Mills gradient flow. It
is interesting to note that the leading $t\to0$~behavior is completely
independent of the detailed definition of the gradient flow; the structure and
coefficients follow solely from the finiteness of the local products and the
renormalizability of the Yang--Mills theory. The sub-leading corrections in the
asymptotic form, i.e., the coefficients $c_1$ and~$c_2$, depend on the detailed
definition of the gradient flow; in the Appendix, we compute the constants
$c_1$ and~$c_2$ and we have
\begin{align}
%   c_1&=\ln\sqrt{\pi}+\frac{7}{16}\simeq1.00986,
   c_1&=\ln\sqrt{\pi}+\frac{7}{22}
   \simeq0.890547,
\label{eq:(4.30)}
\\
%   c_2&=\ln\sqrt{\pi}+\frac{3}{44}-\frac{1}{4}
%   +\frac{b_1}{2b_0^2}\simeq0.812034.
   c_2&=\ln\sqrt{\pi}-\frac{7}{44}
   +\frac{b_1}{2b_0^2}
   \simeq0.834762.
\label{eq:(4.31)}
\end{align}

Finally, a possible method to determine the factor $\ln(\sqrt{8t}\Lambda)$
in~Eq.~\eqref{eq:(4.29)}, i.e., the flow time~$t$ in the unit of the one-loop
$\Lambda$ parameter~\eqref{eq:(4.18)}, for small flow times is to use the
expectation value of the density operator, Eq.~\eqref{eq:(3.6)}. For this
quantity, by applying Eqs.~\eqref{eq:(4.13)} and~\eqref{eq:(4.17)} to the
result of the one-loop calculation, Eqs.~(2.28) and~(2.29)
of~Ref.~\cite{Luscher:2010iy} (specialized to the pure Yang--Mills theory), we
have the asymptotic form,
\begin{equation}
   t^2\left\langle E(t,x)\right\rangle
   \stackrel{t\to0+}{\sim}
   \frac{3(N^2-1)}{128\pi^2}
   \frac{1}{-2b_0\left[\ln(\sqrt{8t}\Lambda)+c\right]},
\label{eq:(4.32)}
\end{equation}
where
\begin{equation}
   c\equiv\ln(2\sqrt{\pi})+\frac{26}{33}-\frac{9}{22}\ln3
   \simeq1.60396.
\label{eq:(4.33)}
\end{equation}
One may use this asymptotic representation for $\ln(\sqrt{8t}\Lambda)$
in~Eq.~\eqref{eq:(4.29)}.\footnote{In practice, one will use
Eq.~\eqref{eq:(4.29)} to compute $t^2\{T_{\mu\nu}\}_R(x)$
from~$t^2U_{\mu\nu}(t,x)$ and~$t^2E(t,x)$. Then, from the value
of~$\sqrt{8t}\Lambda$, one can deduce $\{T_{\mu\nu}\}_R(x)/\Lambda^4$.}
